\documentclass[a4paper,10pt,fleqn]{article}
\usepackage[margin=1in]{geometry}
\usepackage{amssymb}
\usepackage{adjustbox}
\usepackage{natbib}
\usepackage[colorlinks=true,linkcolor=blue,citecolor=blue,urlcolor=blue]{hyperref}
\usepackage{fuzz}
\usepackage{schemapk}
\usepackage{zed-maths}
\usepackage{zed-proof}
\newdimen\savedleftskip
\begin{document}

\section*{theta-Expression}

\noindent The theta expression constructs the binding for a named schema.  When a schema S has components x and y, theta S is the pair (x, y) viewed as a named record.  The primed form theta S' is the after-state binding.

\bigskip

\begin{zed}
[Item]
\end{zed}

\begin{schema}{Store}
items : \power Item
\end{schema}

\noindent A read-only query that returns a snapshot of the current state uses the theta expression to package the before-state as a single output value.

\bigskip

\begin{schema}{SnapshotStore}
\Xi Store \\
snapshot! : Store
\where
snapshot! = \theta Store
\end{schema}

\noindent The Ξ-style frame condition theta $S = theta$ S' asserts that none of the schema components changed.  This is the definition of Xi S expressed using theta directly.

\bigskip

\begin{schema}{StoreFrame}
\Xi Store
\where
\theta Store = \theta Store'
\end{schema}

\end{document}